\section{总结与展望}\label{sec:outlook}
\subsection{经验总结}
\subsubsection{主要成果}
本作品围绕“次日需水预测如何可信地进入灌溉决策”这一问题，完成了从公开数据到可交互决策工具的完整链条：
\begin{enumerate}
\item 基于三地8年ERA5逐日气象构建49维时间安全特征，梯度提升树次日参考蒸散预测MAE为0.7757 mm/日，较昨日值基线降低7.68\%，$R^2$达0.7103，块自助重采样表明改进稳定；
\item 以2023年独立数据完成共形校准，测试期90\%预测区间整体覆盖率91.06\%，并在留出地区检验中保持较好的迁移能力；
\item 将预测区间上界纳入FAO-56水量平衡规则，在180组生长季仿真中实现了缺水风险的稳定下降：基准情景风险日为0，相对墒情阈值策略在全部36个情景中均降低风险；
\item 实现了离线可运行的决策工作台，具备历史回放、参数交互、供水能力检查与决策记录导出功能，配套9项核心测试、9项浏览器测试和61项交付核验。
\end{enumerate}

\subsubsection{关键经验}
\textbf{时间边界是可信评估的前提。}特征窗口“先滞后、后聚合”、按年份切分训练—校准—测试，并用扰动测试自动检查，是避免信息泄漏、使指标可信的基础。

\textbf{分层评价比单一指标更有说服力。}点预测精度、区间覆盖与灌溉结果分层报告，才发现了“点预测略好不一定少缺水，区间上界才稳定降风险”这一对设计最有指导意义的结论。

\textbf{诊断比平均值更能指明方向。}全年91\%的覆盖率看似达标，分月诊断却揭示出夏季覆盖不足，为下一步改进提供了明确目标。

\subsubsection{不足之处}
\begin{table}[H]\centering\small
\caption{当前不足与改进措施}\label{tab:limitations}
\begin{tabularx}{\linewidth}{@{}p{.24\linewidth}XX@{}}\toprule
不足 & 影响 & 改进措施\\\midrule
区间半宽全年固定 & 6—7月覆盖率约78\%\,$\sim$\,79\%，冬季区间偏宽 & 分季节校准、共形分位数回归\cite{cqr}、自适应共形\cite{aci}\\
仅用历史气象输入 & 难以预知天气突变日 & 接入次日数值天气预报作为特征\\
使用再分析数据 & 再分析产品有发布延迟，网格尺度大于地块 & 改用实时预报与地面气象站数据\\
土壤与作物参数为设定值 & 仿真结果依赖参数选择 & 结合墒情传感器与当地农技资料标定\\
日尺度水桶模型 & 不表达分层入渗与日内胁迫 & 引入AquaCrop等作物模型\cite{aquacrop}\\
示范区为三个相邻地区 & 气候类型覆盖有限 & 扩展到西北、东北等不同气候区\\\bottomrule
\end{tabularx}
\end{table}

\subsection{未来规划}
\subsubsection{技术改进}
\textbf{（1）自适应预测区间。}按季节或按预测误差尺度分组校准，或采用共形分位数回归\cite{cqr}让区间宽度随天气条件变化；对存在分布漂移的时间序列，引入自适应共形推断\cite{aci}在线调整覆盖水平，目标是使每个月份的覆盖率都接近名义水平。

\textbf{（2）融合天气预报。}将未来1\,$\sim$\,3天的气温、降水概率、云量和风速预报作为新特征，并在决策时扣减预报降雨，减少降雨当天提前灌溉的情况，同时将预测时效由次日扩展到未来三天，支持轮灌排班。

\textbf{（3）接入田间墒情。}对接土壤水分传感器，用实测数据替代用户输入的亏缺值，并以“预测—观测”偏差在线修正TAW与作物系数，形成“预测—灌溉—观测—修正”的闭环。

\textbf{（4）多地块联合调度。}在共享水源和管网容量约束下，扩展为多地块轮灌优化，按缺水风险排序分配有限的供水能力。

\subsubsection{验证与推广计划}
\begin{table}[H]\centering\small
\caption{分阶段验证与推广计划}\label{tab:roadmap}
\begin{tabularx}{\linewidth}{@{}p{.16\linewidth}XX@{}}\toprule
阶段 & 主要任务 & 阶段目标\\\midrule
模型升级 & 实现自适应区间与天气预报特征，在新的外部年份上重新测试 & 夏季覆盖率接近90\%，天气转折日误差明显下降\\
影子运行 & 选择具备墒情与用水记录的试点地块，系统每日给出建议、农艺人员照常决策 & 积累建议与实际决策的对比数据，验证建议可执行性\\
对照试验 & 在可比地块开展系统建议与常规管理的对照 & 量化实际节水率、产量与能耗变化\\
推广应用 & 面向农技站、合作社和高校课程提供轻量部署版本 & 形成可复制的部署与培训方案\\\bottomrule
\end{tabularx}
\end{table}

\subsubsection{应用前景}
随着高标准农田建设和农业水价综合改革的推进，基层对“按需灌溉”工具的需求日益增长。墒策部署门槛低、计算过程透明，既可作为农技人员面向农户的灌溉参谋，也可嵌入灌区信息化平台，为墒情监测数据提供“下一步该怎么做”的决策出口；在教学场景中，可作为农业水利与智慧农业课程的实践工具。后续团队将按上述计划推进模型升级与试点验证，用田间数据持续检验和完善系统。
